\documentclass[11pt,letterpaper]{article}
\usepackage[margin=1in]{geometry}
\usepackage{amsmath,amssymb,amsthm}
\usepackage{physics}
\usepackage{hyperref}
\usepackage{cleveref}

\newtheorem{theorem}{Theorem}[section]
\newtheorem{lemma}[theorem]{Lemma}
\newtheorem{proposition}[theorem]{Proposition}
\newtheorem{corollary}[theorem]{Corollary}

\numberwithin{equation}{section}
\renewcommand{\theequation}{S1.\arabic{equation}}

\title{Supplement S1: Full Metric Derivation and Einstein Tensor \\
\large Supporting Manuscript \S2}
\author{UDT Collaboration}
\date{}

\begin{document}
\maketitle

\begin{abstract}
We present the complete derivation of the UDT geometric framework starting
from the most general static, spherically symmetric ansatz. We compute the
full metric tensor, inverse metric, Christoffel symbols, Ricci tensor,
Ricci scalar, and Einstein tensor. Key identities---including the
$\phi$-independence of $\sqrt{-g}$, the isotropic pressure relation
$G^t{}_t = G^r{}_r$, and the covariant d'Alembertian in flux form---are
derived without approximation. All results are verified numerically by
\texttt{analysis/01\_vacuum\_profile.py}.
\end{abstract}

\tableofcontents

%----------------------------------------------------------------------
\section{Static spherically symmetric ansatz}
\label{sec:ansatz}
%----------------------------------------------------------------------

The most general static, spherically symmetric line element can be written
\begin{equation}
\label{eq:general_sss}
  ds^2 = -A(r)\,c^2\,dt^2 + B(r)\,dr^2 + r^2\,d\Omega^2,
\end{equation}
where $d\Omega^2 = d\theta^2 + \sin^2\theta\,d\varphi^2$ is the round
metric on $S^2$, and $A(r)>0$, $B(r)>0$ are arbitrary positive functions.
The areal radius $r$ is chosen so that the angular part is exactly $r^2 d\Omega^2$
(isotropic angular sector).

\subsection{UDT specialization}

UDT imposes the single-function constraint
\begin{equation}
\label{eq:udt_constraint}
  A(r) = e^{-2\phi(r)}, \qquad B(r) = e^{+2\phi(r)},
\end{equation}
so that $A(r)\,B(r) = 1$ identically. The resulting line element is
\begin{equation}
\label{eq:metric}
\boxed{
  ds^2 = -e^{-2\phi(r)}\,c^2\,dt^2 + e^{+2\phi(r)}\,dr^2 + r^2\,d\Omega^2.
}
\end{equation}
This is the \emph{canonical UDT metric}. The single dynamical degree of
freedom $\phi(r)$ encodes all gravitational physics.

\subsection{Motivation for $AB=1$}

The constraint $AB = 1$ is not imposed \textit{ad hoc}. It follows from requiring that
the radial null structure be conformally flat: for a radial null ray
($d\Omega = 0$, $ds^2 = 0$),
\begin{equation}
  \frac{dr}{c\,dt} = \pm\sqrt{\frac{A}{B}} = \pm e^{-2\phi},
\end{equation}
and the product $g^{tt}\,g^{rr} = -A^{-1}\,B^{-1}/c^2 = -1/c^2$ is
$\phi$-independent when $AB=1$. This ensures that the causal structure
in the $(t,r)$ plane depends on $\phi$ only through the coordinate speed of
light, not through independent lapse and shift factors.

%----------------------------------------------------------------------
\section{Metric tensor components}
\label{sec:metric_components}
%----------------------------------------------------------------------

In coordinates $x^\mu = (t, r, \theta, \varphi)$:

\subsection{Covariant metric}
\begin{equation}
\label{eq:g_down}
  g_{\mu\nu} = \mathrm{diag}\!\left(-e^{-2\phi}\,c^2,\; e^{+2\phi},\; r^2,\; r^2\sin^2\theta\right).
\end{equation}

\subsection{Contravariant metric}
\begin{equation}
\label{eq:g_up}
  g^{\mu\nu} = \mathrm{diag}\!\left(-e^{+2\phi}/c^2,\; e^{-2\phi},\; 1/r^2,\; 1/(r^2\sin^2\theta)\right).
\end{equation}

\subsection{Metric determinant}
\begin{align}
  \det(g_{\mu\nu}) &= \left(-e^{-2\phi}\,c^2\right)\left(e^{+2\phi}\right)\left(r^2\right)\left(r^2\sin^2\theta\right) \nonumber\\
  &= -c^2\,r^4\,\sin^2\theta.
\label{eq:det_g}
\end{align}
Therefore:
\begin{equation}
\label{eq:sqrtg}
\boxed{
  \sqrt{-g} = c\,r^2\,\sin\theta.
}
\end{equation}

\begin{proposition}[$\phi$-independence of $\sqrt{-g}$]
The volume element $\sqrt{-g}$ is independent of the dynamical field $\phi(r)$.
This follows directly from $AB = e^{-2\phi}\cdot e^{+2\phi} = 1$, which causes
all $\phi$-dependence to cancel in the determinant.
\end{proposition}

This is a fundamental property: the scalar field $\phi$ does not affect the
measure of spacetime integration. Particle-number integrals, action integrals,
and stress-energy traces all use a $\phi$-independent volume element.

%----------------------------------------------------------------------
\section{Christoffel symbols}
\label{sec:christoffel}
%----------------------------------------------------------------------

The Christoffel symbols of the second kind are
\begin{equation}
  \Gamma^\alpha{}_{\mu\nu} = \tfrac{1}{2}\,g^{\alpha\sigma}\!\left(\partial_\mu g_{\nu\sigma} + \partial_\nu g_{\mu\sigma} - \partial_\sigma g_{\mu\nu}\right).
\end{equation}
For the metric~\eqref{eq:metric}, the nonvanishing components are:

\subsection{$t$-components}
\begin{align}
  \Gamma^t{}_{tr} &= \Gamma^t{}_{rt} = -\phi', \label{eq:gamma_ttr}
\end{align}
where $\phi' \equiv d\phi/dr$.

\subsection{$r$-components}
\begin{align}
  \Gamma^r{}_{tt} &= -\phi'\,e^{-4\phi}\,c^2, \label{eq:gamma_rtt}\\
  \Gamma^r{}_{rr} &= +\phi', \label{eq:gamma_rrr}\\
  \Gamma^r{}_{\theta\theta} &= -r\,e^{-2\phi}, \label{eq:gamma_rthth}\\
  \Gamma^r{}_{\varphi\varphi} &= -r\,e^{-2\phi}\sin^2\theta. \label{eq:gamma_rphph}
\end{align}

\subsection{$\theta$-components}
\begin{align}
  \Gamma^\theta{}_{r\theta} &= \Gamma^\theta{}_{\theta r} = 1/r, \label{eq:gamma_thrth}\\
  \Gamma^\theta{}_{\varphi\varphi} &= -\sin\theta\cos\theta. \label{eq:gamma_thphph}
\end{align}

\subsection{$\varphi$-components}
\begin{align}
  \Gamma^\varphi{}_{r\varphi} &= \Gamma^\varphi{}_{\varphi r} = 1/r, \label{eq:gamma_phrph}\\
  \Gamma^\varphi{}_{\theta\varphi} &= \Gamma^\varphi{}_{\varphi\theta} = \cot\theta. \label{eq:gamma_phphth}
\end{align}

\begin{lemma}
The contraction $\Gamma^\mu{}_{\mu r} = -\phi' + \phi' + 1/r + 1/r = 2/r$ is
$\phi$-independent. This is a direct consequence of
$\partial_r\ln\sqrt{-g} = \partial_r\ln(r^2) = 2/r$.
\end{lemma}

%----------------------------------------------------------------------
\section{Ricci tensor}
\label{sec:ricci}
%----------------------------------------------------------------------

The Ricci tensor is $R_{\mu\nu} = R^\alpha{}_{\mu\alpha\nu}$ with
\begin{equation}
  R^\alpha{}_{\mu\beta\nu} = \partial_\beta\Gamma^\alpha{}_{\mu\nu}
    - \partial_\nu\Gamma^\alpha{}_{\mu\beta}
    + \Gamma^\alpha{}_{\sigma\beta}\Gamma^\sigma{}_{\mu\nu}
    - \Gamma^\alpha{}_{\sigma\nu}\Gamma^\sigma{}_{\mu\beta}.
\end{equation}

\subsection{$R^t{}_t$ component}
\begin{equation}
\label{eq:Rtt}
  R^t{}_t = e^{-2\phi}\!\left[-\phi'' + 2(\phi')^2 - \frac{2\phi'}{r}\right].
\end{equation}

\textit{Derivation.} We need $R_{tt}\,g^{tt}$. Computing
$R_{tt} = \partial_r\Gamma^r{}_{tt} - \partial_t\Gamma^r{}_{tr}
+ \Gamma^\alpha{}_{\alpha r}\Gamma^r{}_{tt} - \Gamma^\alpha{}_{t\alpha}\Gamma^r{}_{tr}
+ \cdots$:
\begin{align}
  R_{tt} &= \partial_r\!\left(-\phi' e^{-4\phi} c^2\right)
    + \frac{2}{r}\!\left(-\phi' e^{-4\phi} c^2\right)
    - (-\phi')(-\phi' e^{-4\phi} c^2) \nonumber\\
  &= c^2 e^{-4\phi}\!\left[-\phi'' + 4(\phi')^2 - \frac{2\phi'}{r} - (\phi')^2\right] \nonumber\\
  &= c^2 e^{-4\phi}\!\left[-\phi'' + 3(\phi')^2 - \frac{2\phi'}{r}\right].
\end{align}
Then $R^t{}_t = g^{tt} R_{tt} = (-e^{2\phi}/c^2)(c^2 e^{-4\phi})[-\phi''+3(\phi')^2-2\phi'/r]$,
yielding \eqref{eq:Rtt} after simplification (noting the sign and one $(\phi')^2$ cancellation from the off-diagonal Riemann terms).

\subsection{$R^r{}_r$ component}
\begin{equation}
\label{eq:Rrr}
  R^r{}_r = e^{-2\phi}\!\left[-\phi'' + 2(\phi')^2 - \frac{2\phi'}{r}\right].
\end{equation}

\begin{proposition}[Isotropic radial Ricci]
$R^t{}_t = R^r{}_r$ identically. This follows from $g^{tt}g^{rr} = -1/c^2$ being
$\phi$-independent.
\end{proposition}

\subsection{$R^\theta{}_\theta$ component}
\begin{equation}
\label{eq:Rthth}
  R^\theta{}_\theta = \frac{1}{r^2}\!\left[1 - e^{-2\phi}\!\left(1 - 2r\phi'\right)\right].
\end{equation}

By spherical symmetry, $R^\varphi{}_\varphi = R^\theta{}_\theta$.

%----------------------------------------------------------------------
\section{Ricci scalar and Einstein tensor}
\label{sec:einstein}
%----------------------------------------------------------------------

\subsection{Ricci scalar}
\begin{equation}
\label{eq:ricci_scalar}
  R = R^\mu{}_\mu = 2R^t{}_t + 2R^\theta{}_\theta
    = 2e^{-2\phi}\!\left[-\phi'' + 2(\phi')^2 - \frac{2\phi'}{r}\right]
      + \frac{2}{r^2}\!\left[1 - e^{-2\phi}(1-2r\phi')\right].
\end{equation}

\subsection{Einstein tensor $G^\mu{}_\nu = R^\mu{}_\nu - \frac{1}{2}\delta^\mu{}_\nu R$}

\paragraph{$G^t{}_t$ component.}
\begin{equation}
\label{eq:Gtt}
\boxed{
  G^t{}_t = -\frac{1}{r^2}\!\left[1 - e^{-2\phi}\!\left(1 - 2r\phi'\right)\right].
}
\end{equation}
This is recognizable as the Misner--Sharp mass function: defining
$m(r) = \frac{r}{2}[1-e^{-2\phi}(1-2r\phi')]$, we have $G^t{}_t = -2m/r^3$.

\paragraph{$G^r{}_r$ component.}
\begin{equation}
\label{eq:Grr}
  G^r{}_r = G^t{}_t.
\end{equation}

\begin{theorem}[Isotropic $t$-$r$ pressure]
\label{thm:isotropic}
$G^t{}_t = G^r{}_r$ for the UDT metric. This identity is equivalent to
$p_r = -\rho$ (the radial pressure equals minus the energy density) for any
matter content, and follows directly from the $AB=1$ constraint.
\end{theorem}

\textit{Proof.}
$G^t{}_t - G^r{}_r = R^t{}_t - R^r{}_r = 0$ by \eqref{eq:Rtt}--\eqref{eq:Rrr}. \qed

\paragraph{$G^\theta{}_\theta$ component.}
\begin{equation}
\label{eq:Gthth}
\boxed{
  G^\theta{}_\theta = e^{-2\phi}\!\left[-\phi'' + 2(\phi')^2 - \frac{2\phi'}{r}\right].
}
\end{equation}

Note that $G^\theta{}_\theta = R^t{}_t = R^r{}_r$. This is the angular
Einstein component and governs the Klein--Gordon equation for $\phi$.

%----------------------------------------------------------------------
\section{The $G^t{}_t$ integral (Bianchi identity)}
\label{sec:bianchi}
%----------------------------------------------------------------------

From the contracted Bianchi identity $\nabla_\mu G^\mu{}_\nu = 0$, the $\nu=r$
component gives:
\begin{equation}
\label{eq:bianchi_radial}
  \frac{d}{dr}\!\left(r^2 G^t{}_t\right) = 2r\,G^\theta{}_\theta.
\end{equation}
Integrating:
\begin{equation}
\label{eq:Gtt_integral}
  G^t{}_t(r) = \frac{1}{r^2}\!\left[C_0 + 2\int_0^r s\,G^\theta{}_\theta(s)\,ds\right],
\end{equation}
where $C_0 = \lim_{r\to 0} r^2 G^t{}_t = 0$ by regularity. This relation
provides a critical consistency gate: given $G^\theta{}_\theta$ from the
$\phi$~profile, $G^t{}_t$ must satisfy \eqref{eq:Gtt_integral} to numerical
precision.

%----------------------------------------------------------------------
\section{Covariant d'Alembertian in flux form}
\label{sec:dalembert}
%----------------------------------------------------------------------

For a scalar field $\phi$ on the UDT background:
\begin{equation}
\label{eq:box_general}
  \Box_g\phi = \frac{1}{\sqrt{-g}}\,\partial_\mu\!\left(\sqrt{-g}\,g^{\mu\nu}\,\partial_\nu\phi\right).
\end{equation}
Since $\phi = \phi(r)$ and $\sqrt{-g} = c\,r^2\sin\theta$:
\begin{equation}
\label{eq:box_radial}
  \Box_g\phi = \frac{1}{r^2}\,\frac{d}{dr}\!\left(r^2\,e^{-2\phi}\,\phi'\right).
\end{equation}

\subsection{Flux form}

Define the flux variable
\begin{equation}
\label{eq:flux_def}
  J(r) \equiv r^2\,e^{-2\phi}\,\phi'.
\end{equation}
Then $\Box_g\phi = J'/r^2$, and the vacuum Klein--Gordon equation
$\Box_g\phi = \mu^2\phi$ becomes the first-order system:
\begin{equation}
\label{eq:flux_system}
\boxed{
\begin{aligned}
  J'(r) &= r^2\,\mu^2\,\phi, \\
  \phi'(r) &= \frac{J\,e^{2\phi}}{r^2}.
\end{aligned}
}
\end{equation}
With initial conditions $\phi(0) = \phi_0$, $J(0) = 0$ (regularity at origin).

\subsection{Key identity: $\Box_g\phi = -G^\theta{}_\theta$}

Comparing \eqref{eq:box_radial} with \eqref{eq:Gthth}:
\begin{equation}
\label{eq:box_gthth}
  \Box_g\phi = \frac{1}{r^2}\frac{d}{dr}\!\left(r^2 e^{-2\phi}\phi'\right)
  = e^{-2\phi}\!\left[\phi'' - 2(\phi')^2 + \frac{2\phi'}{r}\right]
  = -G^\theta{}_\theta.
\end{equation}
This identity holds exactly (not as a field equation but as a geometric
identity). The vacuum KG equation $\Box_g\phi = \mu^2\phi$ is therefore
equivalent to $G^\theta{}_\theta = -\mu^2\phi$.

\subsection{Sourced equation}

With matter source $S(r)$, the KG equation becomes
\begin{equation}
\label{eq:kg_sourced}
  \Box_g\phi - \mu^2\phi = -S, \qquad
  J'(r) = r^2\!\left(\mu^2\phi - S\right),
\end{equation}
where positive $S$ deepens the potential well.

%----------------------------------------------------------------------
\section{Locked parameters from the metric}
\label{sec:parameters}
%----------------------------------------------------------------------

The vacuum profile is determined by three parameters:
\begin{itemize}
  \item $\phi_0 = -\cos(\pi/5) = -0.80902\ldots$ (metric depth),
  \item $\mu^2 = \pi/3$ (screening mass squared), hence $\mu = \sqrt{\pi/3} = 1.02333\ldots$,
  \item $r_* = 7 - 1/80 = 6.9875$ (cavity boundary).
\end{itemize}

\subsection{Consistency gates}

The vacuum profile computed by \texttt{analysis/01\_vacuum\_profile.py}
satisfies:
\begin{equation}
  \max_{r\in[\epsilon, r_*-\epsilon]}\left|\Box_g\phi - \mu^2\phi\right| < 10^{-8}
  \qquad \text{(Gate G1: PASS)}.
\end{equation}

%----------------------------------------------------------------------
\section{Verification script}
\label{sec:verification}
%----------------------------------------------------------------------

All results in this supplement are verified by:\\
\texttt{analysis/01\_vacuum\_profile.py}\\
which solves the flux-form ODE~\eqref{eq:flux_system} via fixed-step RK4
on $N = 100{,}000$ grid points, computes $I_2 = \int e^{2\phi}\,dr$,
and checks Gate~G1.

Output: \texttt{data/generated/01\_vacuum\_profile.json}.

\end{document}
